\chapter{Platelet Function Test}

\section*{What Is This Test?}
Platelet-function testing evaluates how well platelets adhere, activate, and aggregate in response to agonists. It can help investigate inherited or acquired bleeding disorders and, in selected situations, assess the antiplatelet effect of medicines such as aspirin or P2Y12 inhibitors.

\section*{Alternative Names}
Platelet Aggregation Study; Platelet Function Screen; PFA-100; PFA-200; P2Y12 Function Test; Platelet Reactivity Test

\section*{How It Is Performed}
Depending on the clinical question, the laboratory may use a platelet-function analyzer, light-transmission aggregometry, whole-blood impedance aggregometry, or a targeted P2Y12 assay. The method and agonists are selected according to the suspected disorder or medication being assessed.

\section*{Why It Is Ordered}
\begin{itemize}
  \item Investigation of mucocutaneous bleeding when the platelet count and routine coagulation tests are normal
  \item Evaluation of suspected von Willebrand disease or an inherited platelet-function disorder
  \item Assessment of platelet inhibition in selected patients receiving antiplatelet therapy
  \item Investigation of unexplained bleeding before specialist-directed procedures when clinically indicated
\end{itemize}

\section*{Preparation, Risks, and Limitations}
The test requires a blood sample and results are highly sensitive to specimen collection, transport, platelet count, hematocrit, anemia, inflammation, and medication exposure. Aspirin, NSAIDs, P2Y12 inhibitors, anticoagulants, and some antidepressants may affect results; medications should not be stopped without clinical direction. A normal screen does not exclude every platelet disorder, and abnormal findings often require confirmatory testing.
\section*{Regulatory Status and Authorization Date}
Regulatory status applies to the specific assay, instrument, manufacturer, and intended use rather than to the test name alone. No single approval date can be assigned to this test category. For a commercial assay, verify the current product in the relevant regulator's database: the U.S. Food and Drug Administration (FDA) clearance, approval, or authorization; India's Central Drugs Standard Control Organisation (CDSCO) licence or permission; the European Union In Vitro Diagnostic Regulation (EU IVDR) CE certificate issued through the applicable conformity-assessment route; or the United Kingdom Medicines and Healthcare products Regulatory Agency (MHRA) registration and UKCA/CE status. Laboratory-developed tests may not have an individual FDA, CDSCO, EU IVDR, or MHRA approval date.
\clearpage
